\section{Esecuzione In-Memory con
Mascheramento}\label{esecuzione-in-memory-con-mascheramento}

Un\textquotesingle altra alternativa per evitare il costo computazionale
e di storage della clonazione dei container consiste
nell\textquotesingle eseguire gli script di verifica direttamente
all\textquotesingle interno dell\textquotesingle ambiente di produzione,
mascherando l\textquotesingle esecuzione su piu fronti, estendendo
l\textquotesingle idea vista in 4.3 per creare
un\textquotesingle ambiente completamente separato per il payload.

\textbf{esecuzione in memoria}: Si mette lo script o il binario in RAM,
utilizzando \passthrough{\lstinline!memfd\_create!}, e lo si esegue
usando \passthrough{\lstinline!fexecve!}, si evitano file persistenti
sul filesystem.

\textbf{mascheramento di \passthrough{\lstinline!/proc!}}: per
nascondere il processo in esecuzione a \passthrough{\lstinline!pspy!} si
inietta nel sistema una libreria a livello utente
\passthrough{\lstinline!/etc/ld.so.preload!} oppure modificando la
variabile d\textquotesingle ambiente
\passthrough{\lstinline!LD\_PRELOAD!}, intercettando le funzioni della
lib \passthrough{\lstinline!C!} che leggono dentro
\passthrough{\lstinline!/proc!}.

\textbf{mascheramento dei metadati del processo}: sovrascrivere il nome
del processo nelle strutture dati del kernel ancora prima che venga
avviato e catturato da \passthrough{\lstinline!pspy!}.

tutta via, queste sono tecniche di oscuramento
dell\textquotesingle esecuzione e possono essere facilmente mascherate e
potrebbero risultare incompatibili con i requisiti
dell\textquotesingle utente:

\begin{itemize}
\tightlist
\item
  l\textquotesingle exploit potrebbe totalmente ignorare
  \passthrough{\lstinline!LD\_PRELOAD!}.
\item
  il file \passthrough{\lstinline!/etc/ld.so.preload!} potrebbe essere
  accessibile allo sfidante, e dunque rimovibile.
\end{itemize}

a livello kernel, esiste comunque un traccia
dell\textquotesingle esecuzione del payload: se l\textquotesingle utente
e\textquotesingle{} in grado di registrare l\textquotesingle esecuzione
delle chiamate di sistema allora puo\textquotesingle{} catturare
l\textquotesingle esecuzione del payload.

Anche se l\textquotesingle eseguibile del verificatore risiede in RAM,
se questo inoltra richiesta HTTP con all\textquotesingle interno il
payload, allora il contenuto puo\textquotesingle{} essere catturato e
letto dallo sfidante.

\begin{itemize}
\tightlist
\item
  il traffico generato puo\textquotesingle{} essere letto da
  \passthrough{\lstinline!tcpdump!}
\item
  oppure la richiesta HTTP puo\textquotesingle{} modificare file
  accessibili allo sfidante.
\end{itemize}

l\textquotesingle esecuzione dell\textquotesingle exploit in ambiente
separato potrebbe comunque portare in crash il servizio.

dunque, l\textquotesingle esecuzione in-memory con mascheramento di
\passthrough{\lstinline!/proc!} si basa su tecniche di evasione fragili
("security through obscurity"). Non offre un vero isolamento: non
protegge i file di log, non isola il traffico di rete e non previene il
crash distruttivo dei servizi live.

La clonazione tramite ZFS CoW in una Dark Sandbox OVN rimane la
soluzione superiore, poiché garantisce un isolamento deterministico a
livello di kernel e hypervisor senza ricorrere a trucchi di evasione in
user-space.
